\documentclass{report}
\usepackage{amsfonts, amsmath}
\newcommand\R{{\mathbb R}}
\newcommand\Q{{\mathbb Q}}
\newcommand\N{{\mathbb N}}
\newcommand\Or{{\mathcal O}}
\newcommand\ve{\varepsilon}
%\pagestyle{empty}
\begin{document}
\large
\centerline{\Huge Fisica Matematica II}
\renewcommand{\thefootnote}{ } 
\renewcommand{\thefootnote}{\arabic{footnote}} 
\vskip.1cm
\centerline{\Large Soluzioni esonero 30-04-2003}
\vskip1cm
\begin{enumerate}
\item La condizione iniziale non \`e una caratteristica. Dunque
\[
\begin{split}
&\dot x=1\\
&\dot y=1-y^2\\
&\dot z=z\\
&x=s ;\quad y=0;\quad z=s.
\end{split}
\]
Dalla seconda segue che 
\[
\frac d{dt}\ln\sqrt{\frac{1+y}{1-y}}=1.
\]
Dunque si ha, per ogni $t\geq 0$,
\[
\begin{split}
&\ln z- x=\ln s-s\\
&\ln z-\ln\sqrt{\frac{1+y}{1-y}}=s.
\end{split}
\]
Che implica
\[
u(x,y)=x\sqrt{\frac{1+y}{1-y}}-\sqrt{\frac{1+y}{1-y}} \ln
\sqrt{\frac{1+y}{1-y}} .
\]
La soluzione \`e unica nel dominio $-1<y<1$.
\item  Se la funzione \`e armonica allora $f''g+g''f=0$ e questo \`e
possibile solo se esiste $\lambda\in\R$ tale che $f''=\lambda f$ e
$g''=-\lambda g$. Le uniche soluzioni di tali equazioni sono
\[
f(x)=\left\{\begin{aligned}
&e^{\sqrt\lambda}x,\; e^{-\sqrt\lambda}x \quad \quad \text{for }
\lambda>0\\
&1,\; x \quad \quad \text{for } \lambda=0\\
&\cos\sqrt{-\lambda} x,\;\sin \sqrt{-\lambda} x\quad \quad \text{for }
\lambda<0, 
	    \end{aligned}\right.
\]
assieme alle analoghe espressioni per $g$. Moltiplicando si ottengo
tutte le possibilit\`a.
\item Il cambio di variabili $(\xi,\eta)=(2x,y)$ trasforma l'equazione
in una equazione di Laplace in un anello. L'esistenza ed unicit\`a
seguono dunque dai soliti argomenti.
\item Sia $u$ una funzione armonica, si consideri $u\circ \Psi$. Un
semplice calcolo mostra che
\[
\partial_{x_i x_i}(u\circ \Psi)=\sum_{k,j}(\partial_{x_k x_j} u)\circ\Psi
A_{ki}A_{ji}.
\]
Dunque si vuole che, per ogni $u$ armonica,
\[
0=\Delta(u\circ\Psi)=\sum_{kj}(A A^*)_{jk}(\partial_{x_k
x_j}u)\circ\Psi=\text{Tr }AA^*U,
\]
Dove $U_{kj}:=(\partial_{x_k x_j}u)\circ\Psi$. Nota che $\text{Tr
}U=0$. 
Ovviamente una possibile soluzione \`e
$AA^*=\lambda\text{Id}$. D'altra parte $u=x_kx_j$, $k\neq j$, sono funzioni
analitiche che implicano che deve essere
\[
0=\sum_{pq}(AA^*)_{pq}\delta_{pk}\delta_{qj}
\]
cio\`e $(AA^*)_{kj}=\lambda_k\delta_{kj}$. Considera ora le funzioni
armoniche $u=x_k^2-x_j^2$. Usando queste si ottine
$\lambda_k=\lambda_j$. Dunque $A^*A$ deve essere un multiplo
dell'identit\`a. Ci\`o implica che $A$ \`e proporzionale ad una
matrice unitaria od ona matrice unitaria per una riflessione.
\end{enumerate}
\end{document}
